\section{Camadas de uma Imagem}

\begin{frame}
  \frametitle{O que são Camadas?}
  Uma imagem é formada por \textbf{camadas imutáveis}: cada uma representa
  um conjunto de modificações no sistema de arquivos.

  \vspace{0.5em}
  Exemplo de uma imagem Python típica:
  \begin{enumerate}
    \item Comandos básicos e gerenciador de pacotes (\texttt{apt})
    \item Runtime Python e \texttt{pip}
    \item Cópia do \texttt{requirements.txt}
    \item Instalação das dependências
    \item Cópia do código-fonte
  \end{enumerate}
\end{frame}

\begin{frame}
  \frametitle{Reutilização de Camadas}
  \begin{block}{Vantagem}
    Se duas aplicações utilizam Python, ambas compartilham as camadas
    de instalação do Python e do \texttt{pip}.
  \end{block}

  \vspace{0.5em}
  \begin{itemize}
    \item Builds mais rápidos
    \item Menos armazenamento necessário
  \end{itemize}
\end{frame}

\begin{frame}
  \frametitle{Empilhamento de Camadas}
  \begin{enumerate}
    \item Cada camada é extraída em seu próprio diretório no sistema de arquivos local
    \item Ao iniciar um container, um \textbf{sistema de arquivos unificado} é criado
          empilhando todas as camadas
    \item O container redireciona seu diretório raiz para essa visão unificada via \texttt{chroot}
  \end{enumerate}
\end{frame}

\begin{frame}
  \frametitle{Criando uma Base Image Manualmente}
  Iniciar container Ubuntu e instalar Node.js:

  \vspace{0.3em}
  \texttt{docker run --name=base-container -ti ubuntu}

  \vspace{0.3em}
  \texttt{apt update \&\& apt install -y nodejs}

  \vspace{0.5em}
  Salvar como nova imagem e verificar camadas:

  \vspace{0.3em}
  \texttt{docker container commit -m "Add node" base-container node-base}

  \vspace{0.3em}
  \texttt{docker image history node-base}

  \vspace{0.5em}
  \begin{block}{Base Image}
    Fundação para construir outras imagens — qualquer imagem pode ser usada como base.
  \end{block}
\end{frame}

\begin{frame}
  \frametitle{Construindo uma App Image}
  A partir da \texttt{node-base}, criar e salvar uma aplicação:

  \vspace{0.3em}
  \texttt{docker run --name=app-container -ti node-base}

  \vspace{0.3em}
  \texttt{echo 'console.log("Hello from an app")' > app.js}

  \vspace{0.5em}
  Salvar com comando padrão definido:

  \vspace{0.3em}
  \texttt{docker container commit -c "CMD node app.js" -m "Add app" app-container sample-app}

  \vspace{0.5em}
  \begin{itemize}
    \item \texttt{-c} define configurações adicionais — aqui, o comando padrão ao iniciar
    \item Na prática, usa-se um \textbf{Dockerfile} para automatizar todos esses passos
  \end{itemize}
\end{frame}